\section{Problem Formulation}
\label{sec:formal}

The key modelling choice is that \emph{source identity is part of the
generative process of the message}, not a label attached to it afterwards:
each source encodes its internal reading through a source-specific
\emph{convention} before transmitting, so what the receiver observes, and
what it means, depends on who sent it and in what context.

\subsection{Generative model}

Let $Y \in \{0,1\}$ be the binary truth with uniform prior, and let there
be $N$ sources. Each source $s$ is described by three parameters, and
keeping them separate is the entire point of the model:

\begin{itemize}
\item A \textbf{convention} $c_s \in \mathcal{C}$: an involutive function
$f_{c_s} : \{0,1\} \times \{0,1\}^d \to \{0,1\}$ mapping the source's
internal reading and an observable context $a \in \{0,1\}^d$ to a message.
Involution means $f_c(f_c(v,a),a) = v$: applying the convention twice
recovers the reading, which is what makes decoding possible at all.
\item A \textbf{stability} $\rho_s \in [0,1]$: the probability that the
source applies its convention on any given transmission rather than
reporting its reading directly.
\item A \textbf{reliability} $r_s \in (1/2, 1)$: the probability that the
source's internal reading equals the truth. This is the only quantity the
trust literature estimates.
\end{itemize}

On each round, source $s$ draws reading $R_s = Y$ with probability
$r_s$ (else $1-Y$), transmits $M_s = f_{c_s}(R_s, a)$ with probability
$\rho_s$ (else $R_s$), and the receiver observes the message $M_s$ together
with the context $a$.

We work with the \textbf{centered message} $X_s = 2M_s - 1 \in \{-1,+1\}$,
so that aggregation is sign-symmetric. Define the log-odds weight
$\lambda_s = \log(r_s/(1-r_s)) > 0$ and, for inversion-only conventions,
the convention sign $\sigma_s = +1$ for honest (identity) conventions and
$\sigma_s = -1$ for inverted conventions.

The full generative process for one round, with source parameters
$\theta = \{(c_s, \rho_s, r_s)\}_{s=1}^N$ fixed, is:
\begin{align}
Y &\sim \mathrm{Bernoulli}(1/2) \\
A_s &\sim P_A \quad \text{independently per source and round} \\
R_s \mid Y & : \quad P(R_s = Y \mid Y) = r_s \\
B_s &\sim \mathrm{Bernoulli}(\rho_s) \\
M_s &= \begin{cases} f_{c_s}(R_s, A_s), & B_s = 1 \\ R_s, & B_s = 0 \end{cases}
\end{align}
The receiver observes $(M_s, A_s)$ for each source. A \textbf{support set}
of $K$ labelled triples $\{(M_s^{(k)}, A_s^{(k)}, Y^{(k)})\}_{k=1}^K$ per
source is available for calibration; test queries are fresh draws from the
same process with the same source parameters.

\textbf{Relation to transition-matrix models.} Fix a source $s$ and a
context $a$, and write the convention in XOR form
$f_{c_s}(v,a) = v \oplus g_{c_s}(a)$ (every convention in our family has
this form; only 8 of the 16 Boolean functions $g$ of two bits occur). The
induced channel $Y \mapsto M_s$ is then a binary symmetric channel with
error rate
\[
\varepsilon_a = \begin{cases} 1 - r_s, & g_{c_s}(a) = 0, \\
(1-\rho_s)(1-r_s) + \rho_s r_s, & g_{c_s}(a) = 1, \end{cases}
\]
since in an unflipped context both branches transmit the reading, while in
a flipped context the direct branch errs with probability $1-r_s$ and the
convention branch with probability $r_s$. Per (source, context) the model
is therefore a strict special case of the per-source, per-context
transition matrices estimated in the crowdsourcing literature
(Section~\ref{sec:related}): an unrestricted model has a free error rate
for each (source, context) pair, while ours ties all contexts through the
two scalars $(\rho_s, r_s)$ and one of eight flip patterns. The signed
effective weight in context $a$ is
$\lambda_a = \log\frac{1-\varepsilon_a}{\varepsilon_a}$: it equals
$+\lambda_s$ in unflipped contexts, $0$ at $\rho_s = 1/2$ in flipped
contexts, and $-\lambda_s$ only at $\rho_s = 1$.
Theorem~\ref{thm:scalar} is a statement about this quantity.

\subsection{Context assumptions}

The context $A_s$ is a $d$-bit vector ($d=2$ in all experiments), drawn
i.i.d.\ from a fixed distribution $P_A$ independently per source, per
round, and independently of the truth $Y$. It is observed without noise,
and the support and test phases share the same $P_A$. No
source chooses or influences the context. These assumptions matter for
Theorem~\ref{thm:scalar}: $P_A$ must put positive mass on the convention's
distinguishing inputs.

\textbf{Edge cases.} With $K = 0$ the posterior of
Section~\ref{sec:method} equals the prior and the predictive is the prior
mixture; nothing in Sections~\ref{sec:theorems}--\ref{sec:method}
requires $K \geq 1$ except the learned calibrators of
Section~\ref{sec:experiments}, which is why all experiments use
$K \geq 1$. Theorem~\ref{thm:consistency}'s predictive claim holds for
$P_A$-almost every query context. The parameter ranges are open
($r_s > 1/2$, $\rho_s > 0$) by design: at $r_s = 1/2$ a source carries no
information, and at $\rho_s = 0$ its convention is never applied, so its
contexts are indistinguishable.

\subsection{Aggregation classes}

We compare two classes of aggregation rules. The \textbf{nonnegative-weight
class}
\begin{equation}
\mathcal{A}_+ = \left\{ \hat{y}_w(X)
= \mathrm{sign}\!\left(\sum_s w_s X_s\right) : w_s \geq 0 \right\}
\end{equation}
captures reliability weighting: each source contributes its (centered)
message scaled by a nonnegative trust weight. The \textbf{signed class}
\begin{equation}
\mathcal{A} = \left\{ \hat{y}_v(X)
= \mathrm{sign}\!\left(\sum_s v_s X_s\right) : v_s \in \mathbb{R} \right\}
\end{equation}
additionally allows negative weights, i.e.\ flipping a source's
contribution. Both classes assign a single scalar per source, independent
of context. In both classes, $\mathrm{sign}(0)$ is resolved by an
independent fair coin, which makes the $\tfrac12$ tie credit in
Lemma~\ref{lem:equalrel} exact.

\textbf{Oracle.} The oracle knows the ground-truth parameters
$(c_s, \rho_s, r_s)$ for every source and applies the Bayes-optimal
aggregation rule directly. No learning, no inference. It is a reference
upper bound, not a competitor.
